\documentclass[12pt]{article}

% ============================================================
% PAQUETES
% ============================================================
\usepackage[margin=2.5cm]{geometry}
\usepackage[utf8]{inputenc}
\usepackage[T1]{fontenc}
\usepackage[spanish]{babel}
\usepackage{enumitem}
\usepackage{setspace}
\usepackage{hyperref}
\usepackage{booktabs}
\usepackage{amsmath}
\usepackage{graphicx}

\setstretch{1.15}

\begin{document}

% ============================================================
% PORTADA
% ============================================================
\begin{titlepage}
    \centering
    
    {\normalsize \textbf{UNIVERSIDAD POPULAR DEL CESAR}}\\[0.3cm]
    {\normalsize \textbf{FACULTAD DE INGENIERÍAS Y TECNOLOGÍAS}}\\[0.3cm]
    {\normalsize \textbf{PROGRAMA DE INGENIERÍA DE SISTEMAS}}\\[2.0cm]
    
    {\large \textbf{SISTEMA DE PRIORIZACIÓN DE BRIGADAS DE SALUD EXTRAMURALES PARA COMUNIDADES INDÍGENAS Y RURALES DISPERSAS}}\\[0.4cm]
    {\normalsize \textit{Documento Técnico de Avance -- Segunda Fase}}\\[2.0cm]
    
    {\normalsize \textbf{PRESENTADO POR:}}\\[0.4cm]
    {\normalsize JORMAN DAVID NORIEGA DÁVILA}\\
    {\normalsize JULIO CESAR RIOS MALDONADO}\\
    {\normalsize GABRIEL OMAR VENEGAS GONZÁLEZ}\\[1.0cm]
    
    {\normalsize \textbf{ASIGNATURA:}}\\
    {\normalsize INTELIGENCIA ARTIFICIAL}\\[0.4cm]
    
    {\normalsize \textbf{GRUPO:}}\\
    {\normalsize 02}\\[0.4cm]
    
    {\normalsize \textbf{DOCENTE:}}\\
    {\normalsize TONNY JIMENEZ}\\[1.5cm]
    
    {\normalsize VALLEDUPAR, CESAR}\\
    {\normalsize 2026}
    
\end{titlepage}

\tableofcontents
\newpage

% ============================================================
% INTRODUCCIÓN
% ============================================================
\section{Introducción}
El acceso equitativo a servicios de salud de atención primaria representa uno de los desafíos más complejos de salud pública en Colombia. En las regiones rurales apartadas y en los territorios colectivos indígenas, la distancia física a los centros urbanos y la precaria infraestructura vial imponen barreras casi insalvables para los pacientes. Como consecuencia, las brigadas médicas extramurales constituyen con frecuencia el único mecanismo de contacto preventivo y curativo con el sistema sanitario institucional.

No obstante, la planeación y el despliegue de estos operativos móviles suelen basarse en criterios empíricos, presiones presupuestales coyunturales o solicitudes informales, lo que ocasiona duplicidad de esfuerzos en ciertas zonas y desatención prolongada en comunidades en extrema vulnerabilidad. El presente documento reporta la consolidación técnica y metodológica de la **Fase 2 (Desarrollo y Modelado)** de un sistema inteligente de priorización y optimización espacial. Mediante la integración de datos abiertos oficiales (RIPS, REPS, Censo DANE 2018 NBI y cartografía DIVIPOLA), se desarrollaron y contrastaron modelos de aprendizaje automático supervisado y un componente prescriptivo de Programación Lineal Entera Mixta (MILP) para transformar la toma de decisiones asistenciales en una política basada en evidencia cuantitativa.

% ============================================================
% ESTADO DEL ARTE
% ============================================================
\section{Estado del Arte}

\subsection{Internacional}
La optimización espacial de clínicas móviles de salud ha cobrado relevancia en la literatura científica reciente. Tanim et al. (2026) desarrollaron un marco de localización-asignación tipo \textit{Maximal Coverage Location Problem} (MCLP) para el despliegue de clínicas móviles en entornos rurales, demostrando un incremento del 90\% en la cobertura poblacional mediante la relocalización sistemática basada en tiempos de viaje. Por su parte, Callaghan et al. (2025) abordaron la planificación de clínicas móviles en la región de Witzenberg (Sudáfrica) implementando un modelo jerárquico de enrutamiento de vehículos con ventanas de tiempo (m-VRPTW) y programación entera para balancear la fatiga del personal médico y minimizar las distancias recorridas.

\subsection{Nacional}
En el contexto latinoamericano y colombiano, Roy y Kar (2022) introdujeron un marco de análisis de decisión multicriterio (MCDA) combinado con análisis de rutas de menor costo para cuantificar las barreras socioeconómicas y geográficas en el acceso a servicios de salud. En Colombia, el uso secundario de las bases de datos masivas del Sistema General de Seguridad Social en Salud (SGSSS), como el Registro de Prestadores (REPS) y los Registros Individuales de Prestación de Servicios (RIPS), ha sido propuesto como insumo estratégico para caracterizar la brecha territorial en zonas ZOMAC y municipios PDET.

\subsection{Local}
En el departamento del Cesar y la región de la Sierra Nevada de Santa Marta y la Serranía del Perijá, la dispersión de comunidades indígenas (pueblos Arhuaco, Wiwa, Kankuamo y Yukpa) plantea retos logísticos severos. La ausencia de vías pavimentadas y la dependencia de transporte fluvial o trochas de herradura demandan herramientas computacionales que optimicen las rutas de salida desde las cabeceras municipales y hospitales de referencia regional como el Hospital Rosario Pumarejo de López en Valledupar.

% ============================================================
% PLANTEAMIENTO DEL PROBLEMA
% ============================================================
\section{Planteamiento del Problema}
La asignación tradicional de brigadas de salud extramurales adolece de tres fallas estructurales:
\begin{enumerate}
    \item \textbf{Inequidad en la cobertura:} Asentamientos con necesidades básicas insatisfechas críticas no reciben visitas por carecer de visibilidad estadística en las secretarías de salud.
    \item \textbf{Ineficiencia operativa y de costos:} Desplazamiento de vehículos y equipos multidisciplinarios a zonas de baja demanda relativa, elevando los costos de combustible y viáticos.
    \item \textbf{Falta de predictibilidad:} Imposibilidad de anticipar picos de demanda ambulatoria según perfiles epidemiológicos y estacionalidad rural.
\end{enumerate}

% ============================================================
% OBJETIVOS
% ============================================================
\section{Objetivos}

\subsection{Objetivo General}
Desarrollar un sistema inteligente de priorización y optimización espacial, basado en modelos de scoring, analítica predictiva de demanda y programación lineal, para planificar la asignación de brigadas de salud extramurales en comunidades indígenas y rurales dispersas.

\subsection{Objetivos Específicos}
\begin{enumerate}
    \item Consolidar y preprocesar las bases de datos oficiales de demanda médica (RIPS), infraestructura (REPS), contexto étnico (MinInterior) e indicadores de Necesidades Básicas Insatisfechas (DANE) a nivel municipal.
    \item Desarrollar un análisis exploratorio de datos (EDA) georreferenciado para identificar patrones de vulnerabilidad y dispersión sanitaria.
    \item Formular, calibrar y evaluar modelos de Machine Learning supervisado para estimar la probabilidad de déficit asistencial y priorización de brigadas.
    \item Diseñar e implementar un modelo de optimización prescriptiva (Programación Lineal Entera Mixta - MILP) que determine el despacho óptimo de brigadas bajo restricciones logísticas y presupuestales.
    \item Construir un prototipo interactivo con soporte de mapas para apoyar la toma de decisiones operativas de las Secretarías de Salud.
\end{enumerate}

% ============================================================
% JUSTIFICACIÓN
% ============================================================
\section{Justificación}
La integración de la Inteligencia Artificial Predictiva (Machine Learning) con la Investigación de Operaciones Prescriptiva permite transitar de un esquema reactivo a un esquema proactivo. Al modelar la vulnerabilidad multidimensional no solo a partir de la cantidad de habitantes, sino incorporando el NBI específico de la zona rural dispersa y las coordenadas geográficas de los municipios, el sistema garantiza una distribución objetiva, verificable y equitativa de los recursos públicos de salud.

% ============================================================
% METODOLOGÍA
% ============================================================
\section{Metodología}

\subsection{Tipo y Método de Investigación}
La investigación es de carácter aplicado, cuantitativo y explicativo, siguiendo las fases de la metodología estándar CRISP-DM adaptada a analítica en salud pública: entendimiento del dominio, ingeniería de datos, modelado, evaluación y despliegue prescriptivo.

\subsection{Ingeniería de Datos y Preprocesamiento}
Se integraron 1.122 municipios de Colombia en un panel anual histórico (2017--2021) acumulando 5.610 observaciones y 74 variables. Las coordenadas GPS decimales fueron normalizadas y el target fue formulado aislando prestadores hub ($\ge 3\times$ la mediana de su estrato) para evitar sesgos urbanos.

\subsection{Partición de Datos 70 / 15 / 15 Agrupada}
Para prevenir fuga de información temporal o autoregresiva, la división se realizó mediante partición agrupada por código DIVIPOLA (\texttt{cod\_mpio}) y estratificada por la prevalencia de prioridad:
\begin{itemize}
    \item \textbf{Entrenamiento (Train):} 785 municipios (70.0\%, 3.925 registros panel).
    \item \textbf{Validación (Val):} 168 municipios (15.0\%, 840 registros panel).
    \item \textbf{Prueba Ciega (Test):} 169 municipios (15.1\%, 845 registros panel).
\end{itemize}

% ============================================================
% RESULTADOS Y MODELADO (FASE 2)
% ============================================================
\section{Desarrollo, Modelado y Evaluación Cuantitativa (Fase 2)}

Se evaluaron 4 modelos representativos con optimización de hiperparámetros vía validación cruzada:
\begin{enumerate}
    \item \textbf{Dummy Classifier:} Estrategia trivial predictora de la clase mayoritaria.
    \item \textbf{Regresión Logística:} Modelo lineal con escalado estándar y balanceo de pesos.
    \item \textbf{Random Forest:} Ensamble de 200 árboles de decisión con hiperparámetros optimizados.
    \item \textbf{HistGradientBoosting:} Modelo de boosting de gradiente de última generación.
\end{enumerate}

\begin{table}[h]
\centering
\small
\caption{Evaluación comparativa de modelos en el conjunto de prueba ciego (Test Set).}
\label{tab:metricas}
\begin{tabular}{lcccccc}
\toprule
\textbf{Modelo Evaluado} & \textbf{Accuracy} & \textbf{Precision} & \textbf{Recall} & \textbf{F1-Score} & \textbf{ROC-AUC} & \textbf{PR-AUC} \\
\midrule
\textbf{HistGradientBoosting} & \textbf{0.7797} & \textbf{0.7413} & 0.5379 & \textbf{0.6234} & \textbf{0.7848} & \textbf{0.7414} \\
Random Forest (Optimizado)    & 0.7283 & 0.6211 & 0.5090 & 0.5595 & 0.7653 & 0.6803 \\
Regresión Logística           & 0.6573 & 0.4952 & \textbf{0.5596} & 0.5254 & 0.6920 & 0.5926 \\
Dummy Baseline (Trivial)      & 0.6610 & 0.0000 & 0.0000 & 0.0000 & 0.5000 & 0.3390 \\
\bottomrule
\end{tabular}
\end{table}

El algoritmo \textbf{HistGradientBoosting} superó a los demás enfoques con un F1-Score de 0.6234 y un ROC-AUC de 0.7848. Destaca su precisión de 74.13\%, garantizando que la gran mayoría de brigadas despachadas llegarán a comunidades genuinamente desatendidas.

\subsection{Componente Prescriptivo: Optimización MILP}
La probabilidad continua derivada del modelo de IA alimentó la función objetivo de un modelo de Programación Lineal Entera Mixta (MILP) resuelto con el solver HiGHS:
\begin{equation}
\max \sum_{i=1}^{N} \left[ P_i \cdot \ln(V_i + 1) - \lambda \cdot d_i \right] x_i
\end{equation}
sujeto a límite de brigadas disponibles ($K=4$) y cobertura mínima de resguardos indígenas ($M=2$). En la prueba piloto para los 25 municipios del Cesar desde Valledupar, el sistema asignó exitosamente las brigadas a Pueblo Bello (Sierra Nevada), Becerril (Perijá), Astrea y Chimichagua.

% ============================================================
% REFERENCIAS
% ============================================================
\section{Referencias Bibliográficas}
\begin{enumerate}[label={[\arabic*]}]
    \item S. H. Tanim, D. L. White, B. Witrick, and L. Rennert, ``Optimizing mobile health clinic placement via geospatial modeling,'' \textit{Public Health in Practice}, vol. 9, p. 100482, 2026.
    \item H. J. Callaghan, L. Potgieter, and N. Le Roux, ``Mobile clinics routing and scheduling in the Witzenberg region of South Africa,'' \textit{PLOS ONE}, vol. 20, no. 3, p. e0284591, 2025.
    \item A. Roy and B. Kar, ``A multicriteria decision analysis framework to measure equitable healthcare access during COVID-19,'' \textit{Journal of Transport \& Health}, vol. 26, p. 101416, 2022.
    \item DANE, ``Censo Nacional de Población y Vivienda (CNPV) 2018: Necesidades Básicas Insatisfechas (NBI) a nivel municipal,'' Departamento Administrativo Nacional de Estadística, Bogotá, 2019.
    \item Ministerio de Salud y Protección Social, ``Registro Especial de Prestadores de Servicios de Salud (REPS) y Sistema RIPS,'' Portal de Datos Abiertos de Colombia, 2021.
\end{enumerate}

\end{document}
